%! TEX root = main.tex

\subsection{구성 요소}

\begin{figure}[ht]
    \centering
    \includegraphics[width=1\linewidth]{figure02.png}
    \caption{System overview}
\end{figure}

본 연구는 영지식 증명 기반의 범위 데이터 검색 시스템을 제안한다. 이 시스템은 발급기관(DI), 데이터 소유자(DP), 데이터 수요자(DC), 검색 서버(SS)로 구성된 서로 다른 역할의 주체들을 포함한다. Figure 1은 본 시스템을 구성하고 있는 주체와 그들 간의 관계를 요약하고 있다. 발급기관이 데이터 소유자에게 데이터를 발급하고 데이터 소유자는 발급기관을 경유해 검색 서버에게 증명 데이터를 전달한다. 그리고 데이터 수요자가 증명에 대한 쿼리를 검색 서버에 전달하면 검색 서버는 그에 상응하는 증명의 목록을 수요자에게 반환하는 구조이다. 즉, 데이터 소유자가 증명을 서버에 제공, 데이터 수요자가 이를 검색하는 구조인 것이다.

\subsubsection{발급기관 (Data Issuer, DI)}
발급기관(DI) 은 사용자에게 데이터를 최초로 발급하는 주체로, 현실 세계에선 공공기관, 병원, 은행 등 인증된 개인정보를 보관하고 또 생성할 수 있는 주체들에 해당한다. 본 연구의 경우, 의료 데이터에 집중하기 때문에 병원 및 검진기관 등으로 고려할 수 있다. DI는 사용자에게 데이터를 발급하면서 본인의 비밀키로 서명을 생성하여 데이터에 첨부한다. 해당 비밀키에 대응되는 공개키는 검색서버(SS)의 데이터베이스에 저장된다. 

\subsubsection{데이터 소유자 (Data Provider, DP)}
데이터 소유자는 데이터에 대한 소유권을 지닌 주체로, DI 로부터 원본을 발급받는다. 현실에선 병원으로부터 의료 데이터를 발급받는 환자에 해당한다. DP는 DI로 부터 데이터를 발급받고 이를 디지털 지갑과 같은 개인화 된 저장소에 보관한다. 이 때, DI는 DP의 승인 하에 데이터를 검색 서버(SS) 에 전달할 수 있고, SS 는 전달받은 영지식 증명을 데이터베이스에 저장한다. 증명 생성 연산은 DI의 시스템에 위임하여 수행하며, DP의 단말기에서 직접 수행하는 것도 가능하다.

\subsubsection{데이터 수요자 (Data Consumer, DC)}
데이터 수요자는 DP의 데이터를 필요로 하는 주체로, 검색을 통해 원하는 속성에 해당하는 데이터를 찾고자 하는 쿼리의 주체이다. 현실에선 의료 데이터를 필요로 하는 제약회사 및 보험회사, 연구기관 등이 이에 해당한다. DC는 검색의 결과로 조회된 증명 데이터들을 로컬에서 직접 검증하여 찾고자 하는 데이터의 존재를 보장받을 수 있다.

\subsubsection{검색 서버 (Search Server, SS)}
검색 서버는 증명을 생성하는 역할과 쿼리를 처리하는 두가지 주된 역할을 맡는다. 우선 DP에게 데이터 원문을 전달받으면 이를 파라미터로 하여 증명 데이터를 생성한다. 생성된 증명 데이터는 검증키 및 메타데이터와 함께 데이터베이스에 저장되며 사용된 원문은 폐기된다. 그리고 DC로부터 속성과 범위를 포함한 쿼리를 수신하면, 그 속성과 범위에 해당하는, 증명 데이터 목록을 데이터베이스로부터 조회하여 반환한다. 

\subsubsection{데이터베이스}

\begin{figure}[ht]
    \centering
    \includegraphics[width=1\linewidth]{figure03-1.png}
    \includegraphics[width=1\linewidth]{figure03-2.png}
    \includegraphics[width=1\linewidth]{figure03-3.png}
    \caption{Database entity structure}
\end{figure}

SS는 검색 시스템을 구현하기 위해 세가지 엔티티 정보들을 데이터베이스에 저장한다. (1) PROOF의 경우, DI로부터 DP의 데이터를 전달받아 생성된 영지식 증명에 대한 정보들을 저장한다. 이는 쿼리의 대상이 되는 데이터이며 원문은 포함하지 않는다. (2) MERKLE의 경우, DI의 공개키로 구성된 머클 트리 정보를 저장하는 테이블이다. 이는 증명하고자 하는 데이터가 시스템에 등록된 DI에 의해 서명된 것인지 익명으로 확인하기 위해 필요한 데이터이다. (3) OWNERSHIP은 특정 PROOF 레코드가 어느 DP의 것인지를 나타내는 소유권 맵핑 정보이다. 해당 정보는 유출될 경우, 특정 사용자(DP)의 데이터 분포가 유추될 가능성이 있기 때문에, 분리된 내부 액세스 전용 데이터베이스로 관리된다. Figure 2는 SS와 연결된 데이터베이스가 가진 세개의 엔티티 구조를 나타내고 있다.

\subsection{데이터 발급 모델}

\begin{figure}[ht]
    \centering
    \includegraphics[width=0.8\linewidth]{figure04.png}
    \caption{Unit data structure}
\end{figure}

본 연구에서 다루는 데이터는 수치 데이터이므로 속성을 나타내는 정보(attribute)와 값을 나타내는 정보(value)가 필수적이다. 또한 의료 정보의 특성상, 기록의 발생 시간이 중요하므로 발급시간에 대한 정보(timestamp)도 필요하다. 그리고 이 세가지 데이터의 무결성을 검증하는 디지털 서명(signature) 역시 필요하다. 따라서 단위 데이터의 원문은 Figure 3와 같은 최소 형식을 따른다고 가정한다. 또한 데이터의 표현 형식은 가장 일반적인 표준인 JSON을 따른다고 가정한다.

\bigskip

예를 들어 ‘\{“attribute”: “attr0001”, “value”: 5.3, “timestamp”: 1768452114, “signature”: “3a8a98…2773aa”\}’와 같이 구성된 데이터가 있다고 가정하자. 그리고 “attr0001”은 “당화혈색소(HbA1c)” 를 참조하는 식별자라고 가정한다. 서명은 “attr0001" + "5.3" + "1768452114” 와 같이 나머지 정보들을 이어 붙여 구성된 메시지와 DI의 비밀키로 생성된다. 추후에 누구든 DI의 공개키를 이용해 이 서명을 검증할 수 있다. 실제로 발급되는 데이터는 이러한 단위 데이터들이 다수 모여 구성된 복합 구조체이며, DP가 가진 디지털 지갑과 같은 안전한 개인 공간에 저장된다.

\begin{figure}[ht]
    \centering
    \includegraphics[width=1\linewidth]{figure05.png}
    \caption{Data issue and proof generation process}
\end{figure}

\bigskip

증명 생성(proof\_gen) 프로세스의 시작에 앞서 ① DI는 DP에게 증명 생성에 대한 동의 여부를 묻는다. 이는 DI가 DP의 데이터를 위임하여 증명을 생성하기 때문이다. 물론 DP가 본인의 단말기에서 직접 증명을 생성해도 문제는 없다. 하지만 DP가 사용하는 단말기는 일반적으로 저사양의 모바일 기기이므로 수십, 수백개의 증명 데이터를 생성하는데 적합한 성능을 기대하기 힘들다. 따라서 데스크탑 및 워크스테이션의 연산능력을 활용할 수 있는 기관인 DI가 이 과정을 위임하는 것을 기본 시나리오로 상정하였다. ②, ③ DP는 자신의 개인키로 요청에 서명하여 이에 동의함을 표시한다. ④ DI는 DP의 데이터에 대해, 임상 참고치를 기반으로 범위 증명 데이터를 생성한다. 여기서 임상 참고치는 SS의 데이터베이스에 공개되어있으며, DI는 이를 불러와 사용한다. ⑤ 생성된 증명 데이터는 SS에게 전달된다. ⑥ SS는 요청을 수신하고 DP의 공개키로 서명을 검증한 후, ⑦ 데이터베이스에 증명 데이터들을 저장한다. Figure 4는 이러한 전체 과정을 도식화 하고 있으며, Figure 5는 DI에서 증명을 생성하고 SS에 전달하는 과정을 의사 코드로 나타내고 있다.

\bigskip

\begin{figure}[ht]
    \centering
    \includegraphics[width=0.8\linewidth]{figure06.png}
    \caption{DI side of proof generation process}
\end{figure}

여기서 미리 정의된 범위 구간은 각 항목이 참조하는 실제 의학적 참고치를 기반으로 한다. 예를 들어 세계보건기구(WHO)의 가이드라인\cite{who2011}에 따르면 당화혈색소의 정상 수치는 6.0\% 미만, 당뇨 전 주의 단계는 6.0 - 6.5\%, 당뇨 진단 기준이 6.5\% 이상이라고 한다. 이를 기준으로 하여 구간 정보를 (0, 6.0), (6.0, 6.5), (6.5, 100)와 같이 미리 정의할 수 있다. 앞서 언급한 사례의 경우, DP의 실제 데이터 수치가 5.3이라면 (0, 6.0)에 대한 범위 증명이 가능하다. 따라서 해당 범위 정보를 파라미터로 하여, 하나의 증명 데이터만 생성하게 될 것이다. 자세한 증명 생성 과정은 4.4에서 설명한다.

\bigskip

이와 같은 방식으로 모든 항목에 대해, 가능한 모든 범위 정보에 대해 증명 생성이 완료되면 이들은 모두 SS에게로 전달된다. 전달된 증명 데이터들은 PROOF 데이터베이스에 (항목, 범위 정보, 증명 데이터)의 단위로 정보를 저장하고, 또한 각 증명 데이터와 이를 제공한 DP의 맵핑 정보를 OWNERSHIP 데이터베이스에 저장한 후, 처리를 마친다.

\subsection{데이터 검색 모델}

\begin{figure}[ht]
    \centering
    \includegraphics[width=1\linewidth]{figure07.png}
    \caption{Proof search process}
\end{figure}

Figure 6는 데이터 검색 모델을 도식화 하고있다. 데이터 검색은 ① 데이터 수요자(DC)가 검색 쿼리를 검색 서버(SS)에 전달하면서 시작된다. 이 쿼리는 찾고자 하는 항목의 식별자(a)와 원하는 구간(r)을 포함한다. ② SS 는 PROOF 데이터베이스에서 쿼리의 항목과 구간에 해당하는 모든 증명 정보를 조회하여 ③ 반환한다. 이때 중요한 점은 증명의 소유자인 DP의 식별자(DP\_ID)는 결과에 포함되지 않는다는 것이다. 즉, 증명의 소유자에 대한 정보는 숨겨진다. DC가 알아야 할 정보는 자신의 쿼리를 만족하는 데이터가 분명히 존재한다는 증거 뿐, 이를 소유한 DP의 정체에 대해 알 필요가 없기 때문이다. Figure 2에서 PROOF 엔티티에 DP\_ID의 속성이 빠진 것은 이런 이유에서 이다. 다만 다중 쿼리의 경우, 다수의 증명에 대한 소유권 입증의 필요가 발생하기 때문에 이를 위한 별도의 OWNERSHIP 테이블을 데이터베이스에서 관리한다. 이에 대한 자세한 내용은 4.5에서 설명한다.

\bigskip

구간 정보는 쿼리에 담긴 구간 정보(r)가 단위 구간 크기보다 클 수 있다. SS는 이를 고려해 쿼리를 만족하는 모든 증명 데이터를 조회해야 한다. 예를 들어 DC가 당화혈색소가 위험구간 이상인 사람의 데이터를 검색하고자 한다면 (6.0, 100)과 같은 구간 정보를 쿼리로 전달할 수 있다. 이는 당뇨 주의 구간과 당뇨 진단 구간을 모두 포함하는 범위이다. 반면 SS의 데이터베이스는 (0, 6.0), (6.0, 6.5), (6.5, 100)의 3개 범위 정보를 취급하고 있다. 이때 SS는 (6.0, 6.5), (6.5, 100) 두개의 범위에 해당하는 PROOF 튜플들을 조회하여 반환해야 한다.

\bigskip

DC는 검색 결과로부터 증명 데이터를 다운받아 로컬 환경에서 zk-SNARK 검증 알고리즘을 수행해, 직접 증명을 검증할 수 있다. 검증에 필요한 컴파일된 회로 바이너리(R1CS) 데이터와 검증키 데이터는 SS에 공개되어있다. 증명 데이터는 데이터 원문이 해당하는 구간 안의 값을 가진다는 사실과 SS의 데이터베이스에 등록된 발급기관(DI)에 의해 서명되었다는 사실을 입증한다. 영지식 증명의 구체적인 구조는 4.4에서 설명한다.

\subsection{범위 증명 구조}

\begin{figure}[ht]
    \centering
    \includegraphics[width=1\linewidth]{figure08.png}
    \caption{Structure of range proof circuit}
\end{figure}

증명 생성 프로세스(proof\_gen)에서 만들어지는 증명($\pi$)은 Figure 7과 같은 구조의 회로(circuit)로부터 생성된다. 이러한 구조의 회로로 증명되는 사실은 (1) 데이터 원문 값이 특정 구간 안에 존재한다는 것과 (2) 데이터 원문에 첨부된 서명이 인증된 발급기관(DI) 의 공개키로 검증되었다는 사실, 두가지이다. 그리고 이 과정에서 DI의 정체를 숨기기 위해 DI의 공개키를 비밀 파라미터로 설정하는 대신, 머클 증명으로 해당 DI의 머클 트리 속 존재를 증명한다.

\bigskip

이를 위해 공개 파라미터로는 구간의 하단 값(range.lower)과 상단 값(range.upper), 데이터의 속성 정보(data.attribute)와 발급 시간(data.timestamp), DI 공개키로 구성된 머클 트리의 루트값(merkel\_root)이 할당된다. 이 정보들은 투명하게 공개되는 값이며, 추후 DC가 검증 시에 직접 확인하여 입력하는 항목들이다. 반면 비밀 파라미터로는 데이터 원문의 값(data.value), 발급 기관(DI)의 공개키(di.public\_key) 그리고 머클 증명에 사용될 머클 경로(merkle\_path)와 머클 비트(merkle\_bits)가 할당된다. 

\bigskip

회로 안에선 먼저 데이터가 주어진 범위의 하단 값과 상단 값의 사이에 존재하는지 확인하는 연산이 수행된다. 이 연산을 통과하면 데이터의 서명을 DI의 공개키로 검증하는 연산이 수행된다. 이 과정에서 사용되는 해시와 서명 알고리즘은 각각 zk-friendly 하다고 알려진 MiMC와 EdDSA를 사용한다. 여기까지의 연산은 데이터의 쿼리 만족 여부와 위변조 여부를 검사하는 계산이다. 다음으로 숨겨진 DI의 공개키가 시스템에 등록된 공개키의 집합에 포함된 것임을 증명하기 위한 머클 증명을 수행한다. 머클 증명은 리프 노드에 위치한 공개키에서 시작해, 주어진 머클 비트(왼쪽/오른쪽 정보)와 경로(형제 노드 정보)를 따라 루트 노드를 향해 점진적인 해시 계산을 수행한다. 최종적으로 계산된 값이 공개 파라미터로 주어진 루트 노드와 동일하다면 DI의 공개키는 머클 트리의 구성원임을 확인할 수 있다. 이렇게 특정 원소의 특정 데이터셋에 대한 포함 여부를 판단하는, 이른바 membership problem을 머클 증명을 통해 해결하면 시간 복잡도를 $O(N)$에서 $O(\log_2{N})$으로 줄이는 효과가 있다.

\bigskip

zk-SNARK 증명을 위해선 증명키와 검증키가 필요하고 이는 시스템 구축 시점에 진행되는 셋업(setup) 과정에서 미리 생성 되었으며, 누구나 공개적으로 접근할 수 있다고 가정한다. 이렇게 미리 생성된 키를 모든 증명/검증 과정에서 반복적으로 사용하는 것은 처리 속도를 높일 뿐 아니라, 악의적인 공격자가 조작된 키로 가짜 증명을 생성하는 것을 방지하는 효과가 있다. 또한 Groth16 프로토콜에서 셋업은 서킷의 구조에 종속되기 때문에 서로 다른 회로를 사용하는 증명들 사이에선 같은 키를 공유할 수 없다. 하지만 본 시스템의 범위 증명은 고정된 구조의 영지식 회로를 사용하기 때문에 단일 셋업을 통해 모든 증명 생성과 검증 작업을 수행할 수 있다.

\bigskip

다만 이러한 시스템을 현실에서 구현할 때는 추가적으로 고려해야 할 사항이 있다. 영지식 증명의 구조가 고정된 머클 트리를 가정하고 있기 때문이다. 머틀 트리는 시스템에 등록된 DI들의 공개키를 저장하는 자료구조이다. 따라서 새로운 DI가 추가되거나 삭제될 때, 혹은 키 유출로 인해 기존의 키를 무효화 하고 새로운 키를 등록해야 할 때 머클 트리의 구조는 지속적으로 변화한다. 이렇게 머클 트리의 구조가 변하면 기존 증명 데이터들의 유효성이 흔들리게 된다. 따라서 새로운 트리 정보가 추가될 때 마다, 상태 정보와 유효기간을 관리하여 새로운 증명과 기존 증명의 검증 유효성을 모두 보장하는 시스템이 필요하다. 또한 키 유출 등으로 더이상 증명의 무결성을 보장할 수 없게 된 경우, 해당 트리의 상태 정보는 폐기된 것으로 간주하며, 이를 통해 만들어진 증명 데이터도 폐기해야 한다. 본 연구는 이를 위해 SS가 관리하는 데이터베이스에 MERKLE 테이블을 추가하여 관리하고 있지만, 이러한 DI 머클 트리 정보에 대한 아카이브를 저장하는 데는 블록체인이 적합할 것으로 보인다. 블록체인은 수정할 수 없는 공개 데이터베이스와도 같기에, 증명의 무결성을 보장하기 위한 정보를 공시하기에 알맞은 수단이기 때문이다. 

\subsection{다중 쿼리의 처리}

다중 쿼리란 하나의 쿼리 안에 다수의 서브 쿼리(속성 – 범위 쌍)가 존재하는 경우를 말한다. 예를 들어 DC는 당뇨병 진단 환자의 데이터를 찾기 위해 당화혈색소가 6.5\% 이상인 동시에 공복혈당이 126 mg/dL이상인 환자의 증명을 검색하려고 할 수 있다. 이 경우 SS는 (당화혈색소 – (6.5, MAX)), (공복혈당 – (126, MAX))의 두 쿼리를 동시에 만족하는 DP의 증명들을 반환해야 한다. 그렇다면 검색 결과는 DP를 기준으로 그룹화된 증명의 목록의 형태를 가질 것이다. 하지만 4.3에서 설명한 것과 같이 DP\_ID는 숨겨지기 때문에 DC는 DP에 대한 그 어떠한 정보도 검색 결과에서 얻을 수 없다. 이로 인해 DC는 그룹으로 묶인 두 범위 증명(당화혈색소와 공복혈당)이 정말로 같은 DP로부터 생성된 것인지에 대한 확신이 필요하다. 악의적인 SS라면 전혀 다른 두 DP의 증명을 제시하며 DC를 기만할 수 있기 때문이다. 이처럼 다중 쿼리를 처리할 때는 '소유권 증명의 문제'가 발생한다.

\begin{figure}[ht]
    \centering
    \includegraphics[width=1\linewidth]{figure09.png}
    \caption{Process of composite query handling}
\end{figure}

\begin{figure}[ht]
    \centering
    \includegraphics[width=1\linewidth]{figure10.png}
    \caption{Commitment for ownership proof}
\end{figure}

이 문제를 해결하기 위해 다시 zk-SNARK 증명이 사용된다. ① 우선 DC가 다중 쿼리를 SS에 전달할 때, 고정된 크기의 임의의 바이트 값인nonce를 함께 전달한다. ② SS는 전달받은 다중 쿼리를 처리하기 위해 데이터베이스에서 PROOF와 OWNERSHIP 테이블을 JOIN 한다. 이렇게 JOIN된 테이블은 DP\_ID 정보를 가지게 된다. 여기서 DC의 쿼리로 필터링 한 튜플들을 DP\_ID를 기준으로 분류한다. ③ 그렇게 튜플들이 DP\_ID로 그룹화 되면, DP\_ID를 숨기기 위한 커밋(commitment)을 계산한다. Figure 9은 이 커밋을 계산하는 과정을 나타내고 있다. 한 그룹에 총 n개의 증명이 존재한다고 할 때, 이 n개 증명들의 식별자 값을 모두 더하여 해시한다(prf\_hash). 그리고 이 해시 값을 DP\_ID, nonce와 더한 후 다시 해시 한다. 이렇게 그룹의 모든 증명에 대한 식별자와 소유자(DP)의 식별자 정보를 함축한 커밋을 계산하되, DC가 제공한 임의의 값(nonce)를 추가하여 이를 랜덤화 하는 것이다. 커밋이 계산되면 튜플들을 분류하고 있던 DP\_ID를 커밋으로 대체한다. ④ 그리고 해당 커밋이 올바르게 계산되었음을 증명하는 영지식 증명을 생성한다. ⑤ 그렇게 커밋과 소유권 증명으로 묶인 범위 증명 항목들은 최종적으로 DC에게 검색 결과로써 반환된다. Figure 8은 다중 쿼리의 처리 과정을 나타내고 있다.

\subsection{소유권 증명 문제}

소유권 증명 문제는 다중 쿼리의 결과로 반환된 임의의 범위 증명 정보들이 동일한 DP\_ID와 연계되어 있음을 증명하는 문제이다. 본 연구는 이를 검색 서버(SS)가 다중 쿼리 결과를 그룹화하고 커밋을 계산할 때, 이 연산이 올바르게 수행되었음을 zk-SNARK를 통해 영지식으로 증명하는 방식으로 해결한다. Figure 10은 여기서 사용되는 산술 회로의 구조를  도식화 하였다.

\begin{figure}[ht]
    \centering
    \includegraphics[width=1 \linewidth]{figure11.png}
    \caption{Structure of ownership proof circuit}
\end{figure}

증명하고자 하는 DP\_ID의 커밋을 $c$, 증명의 식별자들을 ($p_1, \cdots, p_n$) 라고 한다. 공개 파라미터로 $c$와 ($p_1, \cdots, p_n$), 그리고 DC가  쿼리에 포함했던 nonce를 할당한다. 그리고 숨겨야 할 값인 DP\_ID는 비밀 파라미터로 할당한다. 산술 회로는 커밋을 계산하고 이를 파라미터로 주어진 $c$와 비교하는 연산을 수행한다. 두 값이 같을 경우, 결과로 반환된 커밋($c$)은 올바르게 계산된 것임이 증명된다. 참고로 여기서 사용되는 증명키와 검증키 역시 4.4에서의 증명과 마찬가지로 시스템 구축 초기에 진행된 셋업 과정을 통해 미리 생성되었으며, 공개된 것으로 가정한다. 증명이 생성되면 SS는 이를 DC에게 반환하며, DC는 이를 검증하여 커밋이 올바르게 계산되었음을 확인할 수 있다. Figure 11은 이와 같은 과정을 거쳐 최종 생성된 검색 결과 아이템의 한 예시를 보여주고 있다. 실제 검색 결과에는 이와 같은 데이터 구조체가 페이지 크기로 설정된 숫자 이하로 반환된다.

\begin{figure}[ht]
    \centering
    \includegraphics[width=1 \linewidth]{Figure12.png}
    \caption{Example of a single item in a composite query result}
\end{figure}

\subsubsection{소유권 증명 프로세스 최적화 (On-the-fly Batch Proof)}

다중 쿼리를 처리할 때, 실시간으로 소유권 증명을 생성하는 것은 일반적인 구현이라 볼 수 없다. zk-SNARK는 증명 생성의 연산 부담이 큰 기술이기 때문이다. 따라서 SS의 연산 부담을 줄이기 위해 쿼리 결과엔 커밋만을 포함하고 추후 DC의 요구에 따라(on-demand) 증명을 선별적으로 생성하는 것이 더 효율적일 수 있다. 하지만 본 시스템의 소유권 증명 회로는 MiMC와 같은 zk-friendly 해시 연산 위주로 구성된 경량 회로이므로, 수 밀리초(ms) 이내에 증명 생성이 가능하다. 따라서 클라이언트(DC)의 추가적인 요청을 기다리는 대화형(interactive) 방식 대신, SS가 검색 결과를 반환할 때 소유권 증명을 실시간으로 일괄 생성하여 동봉하는 비대화형(non-interactive) 방식으로 최적화하여 통신 지연을 최소화한다. 결과적으로 한 번의 쿼리 요청 및 응답(1 round-trip)만으로 검색과 소유권 검증이 모두 완료되므로, 시스템은 높은 실시간성을 확보할 수 있다.

\subsection{시스템 성능 및 보안성에 대한 고찰}
본 제안 시스템은 기존 연구들과 비교하여 확장성과 실용성 측면에서 다음과 같은 차별점 및 트레이드 오프(Trade-off)를 가진다.

\bigskip

첫째, 압도적인 확장성(Scalability)의 확보이다. 본 시스템의 증명 생성 비용은 전체 데이터베이스 크기(N)가 아닌, 검색 조건을 만족하여 필터링된 결과 그룹의 개수(K)에 선형적으로 비례(O(K))한다. ZKSQL\cite{li2023}과 같은 기존 연구들이 전체 데이터 세트에 대해 O(N) 이상의 다항식 연산을 수행해야 했던 것에 비해 이는 엄청난 성능적 이점이다. 더욱이 각 DP 그룹에 대한 증명 생성은 서로 완전히 독립적이므로, SS 환경에서 병렬 처리를 적용하면 다수의 DP에 대한 소유권 증명 비용을 쉽게 극복하고 대규모 시스템에 적용할 수 있다.

\bigskip

둘째, 실용성과 무결성 간의 현실적인 트레이드 오프이다. 본 시스템에서 증명의 근간이 되는 DP\_ID는 OWNERSHIP이라는 서버 내부의 비공개 데이터베이스에 존재한다. DC 입장에서는 SS가 이 테이블을 조회할 때 정직하게, 모든 튜플을 누락 없이 조회했는지 완벽하게 알 방법은 없다. IntegriDB\cite{zhang2015}나 ZKSQL\cite{li2023} 등의 검증 가능한 데이터베이스 기술을 적용하면 쿼리 결과의 누락이 없음(Completeness)을 수학적으로 보장할 수 있지만, 이는 쿼리당 수 분 이상의 증명 시간을 요구하여 '실시간 검색 서비스'라는 본 연구의 핵심 목적을 훼손한다. 물론 소유권 맵핑과 같은 민감 정보를 데이터베이스에 저장하는 것은 분명한 보안 취약점이다. 이에 대한 구체적인 내용은 5장 실험에서 밝히며, 대응 전략도 간략하게 소개할 것이다.

\bigskip

따라서 본 연구는 SS를 '결과를 임의로 누락할 수는 있지만(Incomplete), 없는 결과를 위조하거나 서로 다른 소유자의 데이터를 섞을 수는 없는(Sound)' 반신뢰(Semi-honest) 주체로 가정하였다. 이를 통해 프라이버시와 위조 불가능성이라는 핵심 보안 요건을 완벽히 달성하면서도 프로토콜을 실시간 처리가 가능할 만큼 경량화 할 수 있었다. 비공개 데이터베이스(OWNERSHIP) 조회의 완전한 무결성까지 실시간 수준으로 보장하는 방법론은 향후 연구 과제로 남긴다.